\Needspace{18\baselineskip}
\section{Transporting replacement sites through isomorphisms}
\label{sec:site-transport}

An address denotes an occurrence, not a permanent numerical edge name. If an
isomorphism exchanges two edges, the selected occurrence must move with it.
This section proves existence of that transport at arbitrary box depth,
congruence when both host and replacement change, and invariance of the exact
freshness condition. Depth bounds and ordered outer interfaces remain fixed.

\subsection{Transport of boundary indices and shallow hosts}

\begin{definition}[Boundary transport]\label{term:boundary-transport}
Given equalities $I=I'$ and $O=O'$ of ordered interfaces, transport a port by
fixing its internal name, or by retaining its numerical input or output index
and using the equality to prove the new bound. Transport a graph to the new
interface by keeping its internal vertices, edges, enumerations, and labels,
and mapping every incidence list through this port transport. This graph is
written $H^{I',O'}$ below.
\end{definition}
All graph invariants follow by eliminating the two interface equalities: the
port map is then the identity, and every invariant is exactly the original
one. Reflexive port transport is pointwise the identity, hence is the identity
function by function extensionality. Reflexive graph transport is equal to the
original graph by record equality and the identity law for list mapping.
Keeping carriers definitionally unchanged makes the construction useful even
before the interface equality proofs have been eliminated.
\lean{Graph/TransportPlug}{8}{12}{Port.cast}
\lean{Graph/TransportPlug}{14}{15}{Port.cast_refl}
\lean{Graph/TransportPlug}{17}{19}{Port.cast_refl_fun}
\lean{Graph/TransportPlug}{22}{39}{Net.reboundary}
\lean{Graph/TransportPlug}{41}{43}{Net.reboundary_refl}

\begin{definition}[Complement restriction]\label{term:complement-restriction}
For a bijection $f:A\to B$ and $a\in A$, restrict $f$ to a bijection
$A\setminus\{a\}\to B\setminus\{f(a)\}$. Its inverse is the restriction of
$f^{-1}$.
\end{definition}
Injectivity shows that a surviving source cannot map to $f(a)$. Conversely, if
$f^{-1}(b)=a$, the right inverse equation gives $b=f(a)$, contradicting survival.
Both inverse laws follow from those of $f$, since subtype elements with equal
underlying values are equal.
\lean{Graph/TransportPlug}{46}{51}{Bijection.except}

\begin{lemma}[Transport of attachment positions]
Let $f:C\cong_R D$ and let $e$ be a selected edge. If its corresponding input
interfaces are equal, mapping the $i$th source port of $e$ by $f$ gives the
$i$th source port of $f(e)$. The analogous assertion holds for targets when the
output interfaces are equal.
\end{lemma}
\begin{proof}
Apply optional list lookup at index $i$ to the source-list equation of $f$.
The typing invariant identifies the source-list length with the input-interface
length, so both lookups are in bounds. Lookup commutes with mapping; cancelling
the optional-result constructor gives the port equation. The interface equality
supplies the bound for the target index. Repeat for the target-list equation.
No disjointness between source and target lists is used.
\end{proof}
\lean{Graph/TransportPlug}{58}{65}{Net.RelIso.sourcePort}
\lean{Graph/TransportPlug}{68}{75}{Net.RelIso.targetPort}

\begin{theorem}[Shallow host congruence]
Suppose $f:C\cong_R D$, $R$ is reflexive, and corresponding selected edges have
equal input and output interfaces. For every replacement $A$ of the source
interface,
\[
 C[e:=A]\cong_R D[f(e):=A^{I',O'}].
\]
\end{theorem}
\begin{proof}
On internal vertices use the disjoint sum of the host vertex bijection and the
identity on replacement vertices. On edges use the complement restriction of
the host edge bijection and the identity on replacement edges. These are the
auxiliary maps \texttt{HostPlug.vertices} and \texttt{HostPlug.edges}.
The whole-graph port map commutes with host inclusion by case analysis on a
port. It also commutes with replacement insertion followed by boundary
transport: internal vertices are immediate, and input and output ports use the
attachment-position lemmas just proved.

For surviving edges, label relatedness is supplied by $f$; for inserted edges,
use reflexivity of $R$. Internal vertex labels are preserved by the host
isomorphism or by identity. For incidence of a surviving edge, rewrite with
$f$'s incidence equation, combine successive list maps, and apply commutation
with host inclusion. For an inserted edge, combine the insertion and whole-port
maps and use commutation with replacement insertion. The same argument proves
the target equation. Self-loops are allowed throughout.
\end{proof}
\lean{Graph/TransportPlug}{83}{84}{Net.HostPlug.vertices}
\lean{Graph/TransportPlug}{86}{88}{Net.HostPlug.edges}
\lean{Graph/TransportPlug}{90}{93}{Net.HostPlug.map_host}
\lean{Graph/TransportPlug}{95}{107}{Net.HostPlug.map_insert}
\lean{Graph/TransportPlug}{112}{161}{Net.RelIso.plugHost}

\Needspace{16\baselineskip}
\subsection{Structural correspondence of addresses}

\begin{definition}[Address interface transport]\label{term:site-interface-transport}
Given $s:\mathsf{Site}(G,r,J,K)$ and equalities $J=J'$, $K=K'$, transport its
interface indices to obtain an address of type $\mathsf{Site}(G,r,J',K')$.
The selected occurrence is unchanged. Plugging $A$ at this transported address
equals plugging the boundary-transported $A$ at the original address.
\end{definition}
The definition is equality elimination. For the plugging equation eliminate
the two equalities and use reflexive graph boundary transport; both sides then
become the same substitution.
\lean{Graph/SiteTransport}{10}{12}{Linear.Site.reinterface}
\lean{Graph/SiteTransport}{15}{18}{Linear.plugAt_reinterface}

\begin{lemma}[Corresponding boxes and their replacement]
If $f:G\cong_{d+1}H$ and edge $e$ of $G$ carries box $(n,m,B)$, then $f(e)$
carries a box $(n,m,B')$ with $B\cong_d B'$. Moreover, replacing these two box
bodies by any $A\cong_d A'$ gives recursively isomorphic whole graphs.
\end{lemma}
\begin{proof}
For existence, inspect the target label in the label-relatedness proof for
$e$. A primitive target contradicts the definition of recursive label
equivalence. For a box target, the constructor of the box relation forces
matching $n,m$ and supplies a body isomorphism. This reasoning constructs an
existential statement, so elimination of the propositional box relation is
legitimate.

For replacement retain the vertex and edge bijections of $f$. At $e$ the new
labels are related by the given isomorphism $A\cong_d A'$. At every other edge
injectivity ensures that its target is different from $f(e)$, so both labels
remain unchanged and their original relation applies. Neither replacement
changes incidence or internal vertex labels; the remaining equations are
exactly those of $f$.
\end{proof}
\lean{Graph/SiteTransport}{22}{36}{Linear.DeepIso.boxAt}
\lean{Graph/SiteTransport}{39}{65}{Linear.DeepIso.replaceBox}

\begin{definition}[Corresponding sites]\label{term:site-correspondence}
The proposition $s\simeq t$, implemented by \texttt{SiteIso}, is generated by
two structural rules. A shallow site selecting $e$ corresponds to the site
selecting $f(e)$ for any recursive host isomorphism $f$; transport the latter's
interface indices back along the signature equalities. For sites inside boxes,
use a recursive outer-host isomorphism, corresponding selected box edges with
equal binding/free metadata, and structurally corresponding inner sites.
Both sites have the same local depth and replacement interface.
\end{definition}
This is a relation on addresses themselves, defined before proving anything
about the results of plugging. An inner correspondence may choose any suitable
isomorphism between the box bodies; the outer relation stores existence of a
body isomorphism, not a distinguished naming map for every inner vertex.
\lean{Graph/SiteTransport}{69}{84}{Linear.SiteIso}

\begin{theorem}[Existence of transported sites]
For every recursive isomorphism $f:G\cong_d H$ and every site
$s:\mathsf{Site}(G,r,J,K)$, there exists
$t:\mathsf{Site}(H,r,J,K)$ with $s\simeq t$.
\end{theorem}
\begin{proof}
Induct on $s$. For a shallow site use the shallow correspondence constructor
and the proved equality of corresponding signatures. For an inside site,
extract the corresponding target box and a body isomorphism using the preceding
lemma. Apply the induction hypothesis to the inner site and that body
isomorphism, then assemble the outer correspondence with the inside constructor.
This follows a finite address down the hierarchy and therefore terminates.
\end{proof}
\lean{Graph/SiteTransport}{87}{94}{Linear.Site.transport_exists}

\begin{definition}[Chosen site transport]\label{term:site-transport}
Write $s^f$ for a site chosen from this existence theorem by classical choice.
It has the original local depth and interface in the target host, and satisfies
$s\simeq s^f$ by the specification theorem for choice.
\end{definition}
The construction is noncomputable. No uniqueness, definitional round-trip law,
or canonical choice of inner box isomorphisms is asserted.
\lean{Graph/SiteTransport}{97}{99}{Linear.Site.transport}
\lean{Graph/SiteTransport}{102}{104}{Linear.Site.transport_spec}

\subsection{Substitution and freshness under transport}

\begin{theorem}[Simultaneous host and replacement congruence]
If $s\simeq t$ and $A\cong_r B$, then
\[
 \mathsf{plugAt}(s,A)\cong_d\mathsf{plugAt}(t,B).
\]
\end{theorem}
\begin{proof}
First take the same replacement on both sides and induct on the structural
correspondence. In the shallow case rewrite plugging at a transported interface
using its equation, then apply shallow host congruence. Recursive label
equivalence supplies the signature equalities and reflexivity required there.
In the inside case the induction hypothesis supplies an isomorphism of the
replaced inner bodies. Lift it through the corresponding outer boxes using
box replacement congruence. This proves \texttt{SiteIso.plug}.
For distinct isomorphic replacements, compose that isomorphism with the
previously proved congruence in the replacement at the fixed target site $t$.
\end{proof}
\lean{Graph/SiteTransport}{107}{117}{Linear.SiteIso.plug}
\lean{Graph/SiteTransport}{120}{125}{Linear.SiteIso.plug_congr}

\begin{lemma}[Quotient substitution agrees after host transport]
For every replacement isomorphism class $[A]$,
$\mathsf{plugAtModuloIso}(s,[A])=
\mathsf{plugAtModuloIso}(s^f,[A])$.
\end{lemma}
\begin{proof}
Use quotient induction to choose a representative $A$. The specification of
$s^f$ gives a structural correspondence; substitution congruence gives a
recursive isomorphism between the two results. Quotient soundness identifies
their classes.
\end{proof}
This is invariance under transporting the host/site pair. A separate quotient
type of host/site pairs has not been constructed. Section~\ref{sec:composition}
proves the nested composition and unit laws.
\lean{Graph/SiteTransport}{128}{132}{Linear.plugAtModuloIso_transport}

\begin{lemma}[Invariance of retained holes]
A recursive host isomorphism preserves and reflects the holes outside a
selected outer edge. Structurally corresponding sites retain exactly the same
hole labels at every depth.
\end{lemma}
\begin{proof}
An outer retained occurrence consists of an edge $x\ne e$ and a hole-observing
position in its label. Map $x$ by the host edge bijection and the position by
the previously constructed label-fibre embedding. Injectivity preserves
$x\ne e$, and the embedding preserves its hole observation. For reflection,
use the inverse edge bijection and the inverse isomorphism's fibre embedding;
the right inverse equation ensures that the preimage survives.

Interface transport does not change retention: eliminate the interface
equalities. Now induct on structural site correspondence. The shallow case
uses this equality elimination and the outer retained-hole equivalence. The
inside case uses the definition of deep retention as a disjunction: either
the occurrence lies under a surviving outer edge, or it is retained by the
inner site. Apply the outer equivalence and the induction hypothesis to the
two disjuncts.
\end{proof}
\lean{Graph/SiteTransport}{135}{147}{Linear.DeepIso.retainedHole_iff}
\lean{Graph/SiteTransport}{150}{153}{Linear.SiteRetainedHole_reinterface}
\lean{Graph/SiteTransport}{156}{164}{Linear.SiteIso.retainedHole_iff}

\begin{lemma}[Invariance of freshness and result validity]
For $s\simeq t$ and $A\cong_r B$, exact freshness holds for $(s,A)$ if and
only if it holds for $(t,B)$. The two substituted results satisfy global
unique-hole validity together, whether or not either original host is valid.
\end{lemma}
\begin{proof}
Freshness forbids a label occurring both in the retained host and in the
replacement. Transport the first occurrence with retained-hole equivalence and
the second with the replacement isomorphism's hole-inventory equivalence. This
proves each direction. For result validity, simultaneous substitution
congruence supplies a recursive isomorphism; apply its preservation and
reflection of unique holes. No original-host validity hypothesis is needed
for this invariance statement.
\end{proof}
\lean{Graph/SiteTransport}{167}{173}{Linear.SiteIso.fresh_iff}
\lean{Graph/SiteTransport}{176}{181}{Linear.SiteIso.plug_unique_iff}

\subsection{A moved address and a wrong-address counterexample}

Use the two-hole chain and its edge swap from the preceding section. The first
hole, formerly edge zero, is edge one after renaming. Define the target deep
site to enter the renamed box and select edge one. The outer boxed isomorphism
and the shallow correspondence for the edge swap prove structural
correspondence with the original deep site.
\lean{Graph/SiteTransportExamples}{9}{11}{Linear.Examples.renamedFirstSite}
\lean{Graph/SiteTransportExamples}{14}{16}{Linear.Examples.firstSites_correspond}
Simultaneous congruence now allows both the host body and the inserted chain to
be renamed independently. The replacement consisting of hole one still
collides with the retained hole one: freshness invariance transports the
previously proved failure of freshness to the new site.
\lean{Graph/SiteTransportExamples}{19}{21}{Linear.Examples.simultaneous_deep_renaming}
\lean{Graph/SiteTransportExamples}{24}{26}{Linear.Examples.renamed_collision_not_fresh}

For contrast within this example, define the wrong target site by keeping the
old numerical name zero. It selects hole one instead of hole zero. After
inserting the one-hole replacement, the wrong result retains hole zero,
witnessed by the surviving edge one inside the renamed box. The exact deep
inventory theorem therefore proves that hole zero is present. In the original
substitution, hole zero is consumed and absent, as already proved. A recursive
isomorphism would preserve hole presence, contradicting these two facts.
Thus the wrong-address substitution is not recursively isomorphic to the
original substitution; moving the address is essential.
\lean{Graph/SiteTransportExamples}{29}{31}{Linear.Examples.wrongRenamedSite}
\lean{Graph/SiteTransportExamples}{35}{42}{Linear.Examples.wrong_address_changes_result}
